\documentclass[11pt]{article}
\usepackage[utf8]{inputenc}
\usepackage[T1]{fontenc}
\usepackage{lmodern}
\usepackage{amsmath,amssymb}
\usepackage{graphicx}
\usepackage{booktabs}
\usepackage{ulem}
\usepackage{fixltx2e}
\usepackage[margin=20mm,a4paper]{geometry}
\usepackage{hyperref}
\hypersetup{colorlinks=true,linkcolor=blue,urlcolor=blue}
\usepackage{kotex}
\title{My Document Converter 샘플 문서}
\author{SHKWON}
\date{2026-09-22}

\begin{document}
\maketitle

\section*{<p>My Document Converter 샘플 문서</p>}\label{pmy-document-converter-샘플-문서p}

\section*{<p>소개 (Introduction)</p>}\label{p소개-introductionp}

이 문서는 \textbf{My Document Converter} 가 지원하는 모든 출력 형식의 샘플을 만드는 데 쓰이는
원본입니다. \emph{강조}, \textbf{굵게}, \sout{취소선}, \verb|인라인 코드|,
H\textsubscript{2}O 와 x\textsuperscript{2}, 그리고
\href{https://example.com}{링크}를 담고 있습니다. This paragraph is in English so
that Latin-only formats still show something readable.

\subsection*{<p>목록 (Lists)</p>}\label{p목록-listsp}

• 첫 번째 항목

• 두 번째 항목 • 중첩된 항목 • 또 하나의 중첩 항목

• 세 번째 항목

1. 순서가 있는 항목

2. 두 번째

3. 세 번째

• 끝난 일

• 남은 일

\textbf{용어} 용어에 대한 정의입니다.

\textbf{Pandoc} 범용 문서 변환기. 이 프로그램은 같은 형식 집합을 자체 엔진으로 다룹니다.

\subsection*{<p>코드 (Code)</p>}\label{p코드-codep}

\verb|function greet(name) {|

\verb|  console.log(`Hello, ${name}!`)|

\verb|}|

\verb|greet('world')|

\subsection*{<p>표 (Table)</p>}\label{p표-tablep}

\begin{table}[htbp]
\centering
\begin{tabular}{llll}
\toprule
형식 & 확장자 & 읽기 & 쓰기 \\
\midrule
Markdown & .md & O & O \\
HTML & .html & O & O \\
DOCX & .docx & O & O \\
EPUB & .epub & O & O \\
\bottomrule
\end{tabular}
\end{table}

\subsection*{<p>인용과 수식 (Quote and math)</p>}\label{p인용과-수식-quote-and-mathp}

\begin{quote}
좋은 문서는 어떤 형식으로도 읽힙니다. --- 누군가
\end{quote}

인라인 수식 \verb|E = mc^2| 과 블록 수식:

\verb|\int_0^1 x^2\,dx = \frac{1}{3}|

\subsection*{<p>각주와 그림 (Footnote and figure)</p>}\label{p각주와-그림-footnote-and-figurep}

각주가 붙은 문장입니다.\footnote{[1]} 자동 링크 \url{https://pandoc.org} 와
\url{mailto:knix008@naver.com} 도 있습니다.

\includegraphics[width=\linewidth,keepaspectratio]{../public/app-icon.svg}

마지막 문단입니다. 줄 끝에 두 칸을 두면 강제 줄바꿈이 됩니다.

\end{document}
